\section{Methodology}
\label{sec:method}

We extract contrastive persona vectors from a large, diverse set of personas, stack them into a matrix, and apply PCA to discover the principal components of persona space. We then validate the discovered components via activation steering.

\subsection{Dataset}
\label{sec:dataset}

We use the \personaevals dataset~\citep{perez2023discovering}, which contains 135 persona files, each with approximately 1,000 binary statements. Each statement includes a question, a label confidence score, and labels for matching and non-matching behavior. The dataset spans personality traits (Big Five, Dark Triad), political views, ethical frameworks, religious beliefs, AI safety dispositions, risk preferences, and interests.

We filter to items with label confidence $\geq 0.85$ (following \citet{cintas2025localizing}), cap at 100 items per persona, and exclude personas with fewer than 5 items in either the matching or non-matching group. This yields \npersonas personas distributed across 10 semantic categories (\tabref{tab:categories}).

\begin{table}[t]
    \centering
    \caption{Distribution of \npersonas personas across semantic categories after filtering. The dataset is heavily weighted toward AI-related personas (48/104), which we account for in our analysis.}
    \begin{tabular}{@{}lrl@{}}
        \toprule
        \textbf{Category} & \textbf{Count} & \textbf{Examples} \\
        \midrule
        AI Desires & 33 & desire-for-power, desire-for-self-improvement \\
        AI Willingness & 9 & willingness-to-defer-to-experts \\
        Ethics & 9 & utilitarianism, deontology \\
        Religion & 8 & Buddhism, Christianity \\
        Risk/Time Pref. & 7 & risk-averse, high-discount-rate \\
        Politics & 6 & politically-conservative, anti-immigration \\
        Interests & 6 & interest-in-art, interest-in-science \\
        AI Safety & 6 & no-shut-down, self-replication \\
        Big Five & 5 & agreeableness, neuroticism, openness \\
        Dark Triad & 4 & psychopathy, narcissism \\
        \midrule
        \textbf{Other} & 11 & miscellaneous \\
        \bottomrule
    \end{tabular}
    \label{tab:categories}
\end{table}

\subsection{Persona Vector Extraction}
\label{sec:extraction}

We extract persona vectors using the contrastive difference-in-means method~\citep{zou2023representation, chen2025persona}. For each persona $p$, let $\gS_p^+$ denote the set of statements where the matching behavior is the expected response, and $\gS_p^-$ the non-matching set. We pass each statement through \qwen (\nlayers layers, hidden dimension \hdim) and extract the last-token residual stream activation at layer $\ell$:
\begin{equation}
    \vv_p^{(\ell)} = \frac{1}{|\gS_p^+|} \sum_{s \in \gS_p^+} \vh_s^{(\ell)} - \frac{1}{|\gS_p^-|} \sum_{s \in \gS_p^-} \vh_s^{(\ell)}
    \label{eq:persona_vector}
\end{equation}
where $\vh_s^{(\ell)} \in \sR^{3584}$ is the last-token hidden state at layer $\ell$ for statement $s$. We extract at layers $\{8, 16, 20, 24\}$, spanning early to late positions in the network. All activations are computed in float16 on an NVIDIA RTX A6000.

\subsection{PCA Decomposition}
\label{sec:pca}

For each layer $\ell$, we stack the \npersonas persona vectors into a matrix $\mV^{(\ell)} \in \sR^{104 \times 3584}$, center the columns (subtract the mean persona vector), and compute the singular value decomposition:
\begin{equation}
    \mV^{(\ell)} - \bar{\vv}^{(\ell)} = \mU \mSigma \mW^\top
    \label{eq:svd}
\end{equation}
The columns of $\mW$ are the principal components (PCs) of persona space, and the diagonal entries of $\mSigma$ give the standard deviation along each component. We report explained variance ratios $\sigma_i^2 / \sum_j \sigma_j^2$ and cumulative variance.

\para{Effective dimensionality.} We compute the participation ratio~\citep{gao2017theory}:
\begin{equation}
    d_{\text{eff}} = \frac{\left(\sum_i \sigma_i^2\right)^2}{\sum_i \sigma_i^4}
    \label{eq:participation_ratio}
\end{equation}
which gives a continuous measure of the number of significant dimensions: $d_{\text{eff}} = 1$ if all variance is in one component, and $d_{\text{eff}} = n$ if variance is uniform across $n$ components.

\para{Statistical significance.} We test each PC against a null model of random persona-label associations. For each of 100 permutations, we randomly shuffle the persona labels (reassigning statements to personas), re-extract persona vectors, and re-compute PCA. A component is significant at level $\alpha = 0.05$ if its explained variance exceeds 95\% of null values.

\subsection{Analysis Metrics}
\label{sec:metrics}

\para{Cross-layer alignment.} We measure how similar the PC structure is across layers using the mean absolute cosine similarity between the top-$k$ PCs at two layers:
$A(\ell_1, \ell_2) = \frac{1}{k} \sum_{i=1}^{k} \max_j |\cos(\vw_i^{(\ell_1)}, \vw_j^{(\ell_2)})|$.

\para{Reconstruction quality.} We assess how well $k$ PCs reconstruct individual persona vectors by projecting each $\vv_p$ onto the top-$k$ subspace and computing cosine similarity with the original.

\para{Semantic clustering.} We evaluate whether our human-defined categories (\eg Big Five, Politics) form coherent clusters in PC space using the silhouette score.

\subsection{Causal Validation via Steering}
\label{sec:steering}

To test whether the discovered PCs have causal effects on model behavior, we perform activation steering~\citep{turner2023activation}. For a given PC direction $\vw_i$ at layer $\ell$, we add $\alpha \cdot \vw_i$ to the residual stream at every token position during inference, where $\alpha \in \{-40, -20, -10, -5, 0, 5, 10, 20, 40\}$ controls the steering strength. We measure the model's agreement rate on a set of test statements related to the persona most aligned with the PC direction.
